\ifdefined\gkver\else\def\gkver{student}\fi
\documentclass[\gkver]{gaokaozhenti}
\begin{document}

\chapter{2009年上海理综}
%% paperType: 理综物理部分

\begin{enumerate}
\item
%% number: 4
%% paperName: 2009年上海理综第 4 题 3 分
%% typeId: 1
%% type: 单选题
%% chapter: 宇宙恒星演化
%% point: 天体的演化
%% method: 无
%% score: 3
%% degree: 850
%% duplicateId: 0
%% body:
一颗恒星的寿命取决于它的质量，其关系如右图所示。根据图中信息可以判断\xzanswer{C}

①质量大的恒星寿命较短

②质量小的恒星寿命较短

③太阳的寿命约为10亿年

④太阳的寿命约为100亿年

\onepicture[w=4.5cm, align=right]{figs/04.png}

\fourchoices[answer=C]
{①③}
{②③}
{①④}
{②④}

%% answer: C
%% memo:
\memoanswer{
本题原卷及材料中未提供详解，待补充。
}

\item
%% number: 5
%% paperName: 2009年上海理综第 5 题 3 分
%% typeId: 1
%% type: 单选题
%% chapter: 电路
%% point: 逻辑电路
%% method: 无
%% score: 3
%% degree: 850
%% duplicateId: 0
%% body:
为了保障行驶安全，一种新型双门电动公交车安装了如下控制装置：只要有一扇门没有关紧，汽车就不能启动。如果规定：车门关紧时为“1”，未关紧时为“0”；当输出信号为“1”时，汽车可以正常启动行驶，当输出信号为“0”时，汽车不能启动。能正确表示该控制装置工作原理的逻辑门是\xzanswer{A}

\fourchoices[answer=A]
{与门}
{或门}
{非门}
{与非门}

%% answer: A
%% memo:
\memoanswer{
本题原卷及材料中未提供详解，待补充。
}

\item
%% number: 6
%% paperName: 2009年上海理综第 6 题 3 分
%% typeId: 1
%% type: 单选题
%% chapter: 光学
%% point: 光的电磁说
%% method: 无
%% score: 3
%% degree: 650
%% duplicateId: 0
%% body:
英国科学家瑞利于 1871 年证明：一束光穿过大气 $x$ 距离后，其强度从 $I_{0}$ 下降为 $I(x)$ 的公式为 $I(x)=I_{0}e^{-\alpha x}$，其中 $\alpha=\frac{2\omega^{4}}{3\pi c^{4}N}|n-1|^{2}$ 叫做吸收系数，式中 $\omega$ 为光的频率，$c$ 为光速，标准状况下，$N=2.69\times10^{19}$ 个/$\Ucmc$，$n-1=2.78\times10^{-4}$。定义 $A=\alpha^{-1}$，叫做衰减长度，它表示光经过 $A$ 距离后其强度降低到原来的 $\frac{1}{e}=0.368$。根据以上信息，结合所学知识可以判断\xzanswer{B}

\fourchoices[answer=B]
{可见光中衰减最厉害的是红光}
{可见光中衰减最厉害的是紫光}
{可见光中衰减最厉害的是黄绿光}
{不同颜色的光衰减程序基本相同}

%% answer: B
%% memo:
\memoanswer{
本题原卷及材料中未提供详解，待补充。
}

\item
%% number: 7
%% paperName: 2009年上海理综第 7 题 3 分
%% typeId: 1
%% type: 单选题
%% chapter: 运动和力
%% point: 变加速直线运动
%% method: 无
%% score: 3
%% degree: 550
%% duplicateId: 0
%% body:
右图为蹦极运动的示意图。弹性绳的一端固定在 O 点，另一端和运动员相连。运动员从 O 点自由下落，至 B 点弹性绳自然伸直，经过合力为零的 C 点到达最低点 D，然后弹起。整个过程中忽略空气阻力。分析这一过程，下列表述正确的是\xzanswer{B}

①经过 B 点时，运动员的速率最大

②经过 C 点时，运动员的速率最大

③从 C 点到 D 点，运动员的加速度增大

④从 C 点到 D 点，运动员的加速度不变

\onepicture[w=1.6cm, align=right]{figs/07.png}

\fourchoices[answer=B]
{①③}
{②③}
{①④}
{②④}

%% answer: B
%% memo:
\memoanswer{
本题原卷及材料中未提供详解，待补充。
}

\item
%% number: 8
%% paperName: 2009年上海理综第 8 题 3 分
%% typeId: 1
%% type: 单选题
%% chapter: 直线运动
%% point: DIS测瞬时速度
%% method: 无
%% score: 3
%% degree: 650
%% duplicateId: 0
%% body:
小明同学在学习了 DIS 实验后，设计了一个测物体瞬时速度的实验，其装置如下图所示。在小车上固定挡光片，使挡光片的前端与车头齐平、将光电门传感器固定在轨道侧面，垫高轨道的一端。小明同学将小车从该端同一位置由静止释放，获得了如下几组实验数据。

\begin{table}[htbp]
\centering
\begin{tblr}{|c|c|c|c|}
\hline
实验次数 & 不同的挡光片 & 通过光电门的时间（$\Us$） & 速度（$\Ums$） \\
\hline
第一次 & Ⅰ & 0.23044 & 0.347 \\
\hline
第二次 & Ⅱ & 0.17464 & 0.344 \\
\hline
第三次 & Ⅲ & 0.11662 & 0.343 \\
\hline
第四次 & Ⅳ & 0.05850 & 0.342 \\
\hline
\end{tblr}
\end{table}

\onepicture[w=6.5cm, align=right]{figs/08.png}

则以下表述正确的是\xzanswer{D}

①四个挡光片中，挡光片Ⅰ的宽度最小

②四个挡光片中，挡光片Ⅳ的宽度最小

③四次实验中，第一次实验测得的速度最接近小车车头到达光电门时的瞬时速度

④四次实验中，第四次实验测得的速度最接近小车车头到达光电门时的瞬时速度

\fourchoices[answer=D]
{①③}
{②③}
{①④}
{②④}

%% answer: D
%% memo:
\memoanswer{
本题原卷及材料中未提供详解，待补充。
}

\item
%% number: 9
%% paperName: 2009年上海理综第 9 题 3 分
%% typeId: 1
%% type: 单选题
%% chapter: 电磁感应
%% point: 右手定则
%% method: 无
%% score: 3
%% degree: 550
%% duplicateId: 0
%% body:
信用卡的磁条中有一个个连续的相反极性的磁化区，每个磁化区代表了二进制数 1 或 0，用以储存信息。刷卡时，当磁条以某一速度拉过信用卡阅读器的检测头时，在检测头的线圈中会产生变化的电压（如图 1 所示）。当信用卡磁条按图 2 所示方向以该速度拉过阅读检测头时，在线圈中产生的电压随时间的变化关系正确的是\xzanswer{B}

\onepicture[w=8.5cm]{figs/09a.png}

\fourchoices[answer=B, ispicture=true, h=1.6cm]
{figs/09b.png}
{figs/09c.png}
{figs/09d.png}
{figs/09e.png}

%% answer: B
%% memo:
\memoanswer{
本题原卷及材料中未提供详解，待补充。
}

\item
%% number: 43
%% paperName: 2009年上海理综第 43 题 3 分
%% typeId: 1
%% type: 单选题
%% chapter: 曲线运动
%% point: 线速度角速度
%% method: 无
%% score: 3
%% degree: 750
%% duplicateId: 0
%% body:
右图为一种早期的自行车，这种不带链条传动的自行车前轮的直径很大，这样的设计在当时主要是为了\xzanswer{A}

\onepicture[w=4.0cm, align=right]{figs/43.png}

\fourchoices[answer=A]
{提高速度}
{提高稳定性}
{骑行方便}
{减小阻力}

%% answer: A
%% memo:
\memoanswer{
本题原卷及材料中未提供详解，待补充。
}

\item
%% number: 44
%% paperName: 2009年上海理综第 44 题 4 分
%% typeId: 3
%% type: 填空题
%% chapter: 曲线运动
%% point: 线速度角速度
%% method: 无
%% score: 4
%% degree: 750
%% duplicateId: 0
%% body:
自行车的设计蕴含了许多物理知识，利用所学知识完成下表。

\begin{table}[htbp]
\centering
\begin{tblr}{|c|c|}
\hline
自行车的设计 & 目的（从物理知识角度） \\
\hline
车架用铝合金、钛合金代替钢架 & 减轻车重 \\
\hline
车胎变宽 & \tkanswer[2.5cm]{} \\
\hline
自行车后轮外胎上的花纹 & \tkanswer[2.5cm]{} \\
\hline
\end{tblr}
\end{table}

%% answer:
\jdanswer{
减小压强（提高稳定性）；增大摩擦（防止打滑；排水）
}

%% memo:
\memoanswer{
本题原卷及材料中未提供详解，待补充。
}

\item
%% number: 45
%% paperName: 2009年上海理综第 45 题 6 分
%% typeId: 3
%% type: 填空题
%% chapter: 曲线运动
%% point: 线速度角速度
%% method: 无
%% score: 6
%% degree: 450
%% duplicateId: 0
%% body:
小明同学在学习了圆周运动的知识后，设计了一个课题，名称为：快速测量自行车的骑行速度。他的设想是：通过计算踏脚板转动的角速度，推算自行车的骑行速度。经过骑行，他得到如下的数据：在时间 $t$ 内踏脚板转动的圈数为 $N$，那么脚踏板转动的角速度 $\omega=\tkanswer[2cm]{}$；要推算自行车的骑行速度，还需要测量的物理量有\tkanswer[4cm]{}；自行车骑行速度的计算公式 $v=\tkanswer[3cm]{}$。

\onepicture[w=9cm, align=right]{figs/45.jpg}

%% answer:
\jdanswer{
$2\pi\frac{N}{t}$；牙盘的齿轮数 $m$、飞轮的齿轮数 $n$、自行车后轮的半径 $R$（牙盘的半径 $r_{1}$、飞轮的半径 $r_{2}$、自行车后轮的半径 $R$）；$\frac{m}{n}R\omega$ 或 $2\pi\frac{mN}{nt}R$（$2\pi\frac{r_{1}N}{r_{2}t}R$ 或 $\frac{r_{1}}{r_{2}}R\omega$）
}

%% memo:
\memoanswer{
本题原卷及材料中未提供详解，待补充。
}

\item
%% number: 46
%% paperName: 2009年上海理综第 46 题 4 分
%% typeId: 3
%% type: 填空题
%% chapter: 功和能
%% point: 交通工具启动
%% method: 无
%% score: 4
%% degree: 650
%% duplicateId: 0
%% body:
与普通自行车相比，电动自行车骑行更省力。下表为某一品牌电动自行车的部分技术参数。

\begin{table}[htbp]
\centering
\begin{tblr}{|c|c|c|c|}
\hline
\SetCell[c=2]{c} 规格 & & \SetCell[c=2]{c} 后轮驱动直流永磁铁电机 & \\
\hline
车型 & 14 电动自行车 & 额定输出功率 & $200\UW$ \\
\hline
整车质量 & $40\Ukg$ & 额定电压 & $48\UV$ \\
\hline
最大载重 & $120\Ukg$ & 额定电流 & $4.5\UA$ \\
\hline
\end{tblr}
\end{table}

在额定输出功率不变的情况下，质量为 $60\Ukg$ 的人骑着此自行车沿平直公路行驶，所受阻力恒为车和人总重的 0.04 倍。当此电动车达到最大速度时，牵引力为 \tkanswer[1.5cm]{} $\UN$，当车速为 $2\Ums$ 时，其加速度为 \tkanswer[1.5cm]{} $\Umsq$（$g=10\Umsq$）。

%% answer:
\jdanswer{
40；0.6
}

%% memo:
\memoanswer{
本题原卷及材料中未提供详解，待补充。
}

\item
%% number: 47
%% paperName: 2009年上海理综第 47 题 3 分
%% typeId: 1
%% type: 单选题
%% chapter: 气体性质
%% point: 能源的开发与利用
%% method: 无
%% score: 3
%% degree: 850
%% duplicateId: 0
%% body:
以自行车代替汽车出行，可以减少我们现代生活中留下的“碳足迹”，积极应对全球气候变暖的严峻挑战。我们的各种行为留下的“碳足迹”可以用直观的“碳足迹计数器”进行估算。比如：开车的二氧化碳排放量（$\Ukg$）$=$汽油消耗升数$\times 2.2$。设骑车代替开车出行 $100\Ukm$，可以节约 $9\UL$，则可以减排的二氧化碳约\xzanswer{B}

\fourchoices[answer=B]
{$100\Ukg$}
{$20\Ukg$}
{$9\Ukg$}
{$2.2\Ukg$}

%% answer: B
%% memo:
\memoanswer{
本题原卷及材料中未提供详解，待补充。
}

\end{enumerate}

\end{document}
